\begin{figure}[t]
\centering
\begin{tikzpicture}[scale=1.0]
  % Define colors (brand colors from spec)
  \colorlet{garnet}{Garnet}
  \colorlet{rose}{Rose}
  \colorlet{seablue}{Atlantic}
  \colorlet{neutral70}{black!70}
  \colorlet{neutral50}{black!50}

  % ============ LEFT SIDE: RADAR SIDE VIEW ============

  % Title for left section
  \node[anchor=south, font=\bfseries\Large] at (3.5, 9.5) {Radar Side View (Elevation Plane)};

  % Sea surface (horizontal line)
  \draw[line width=1.5pt, seablue] (0, 1.5) -- (7, 1.5);
  \node[anchor=north, font=\small, seablue] at (7.3, 1.5) {Sea Surface};

  % Radar antenna on left
  \draw[line width=2pt, garnet] (0.5, 1.5) circle (0.15);
  \draw[line width=1.5pt, garnet] (0.5, 1.5) -- (0.5, 0.8);
  \node[anchor=north, font=\small, garnet] at (0.5, 0.5) {Radar};

  % Radar fan beam (wide beamwidth in elevation)
  % Upper beam edge
  \draw[line width=1.5pt, garnet, densely dashed] (0.5, 1.5) -- (6, 4.2);
  % Lower beam edge
  \draw[line width=1.5pt, garnet, densely dashed] (0.5, 1.5) -- (6, 0.2);

  % Filled beam (semi-transparent)
  \fill[garnet, opacity=0.15] (0.5, 1.5) -- (6, 4.2) -- (6, 0.2) -- cycle;

  % Beam angle annotations
  \node[anchor=east, font=\small, garnet] at (0.2, 3.0) {$\pm 12.5^{\circ}$};
  \draw[garnet, line width=0.5pt] (0.35, 2.8) -- (0.35, 3.2);

  % True ship target (at sea surface)
  \draw[line width=1.5pt, neutral70] (4.5, 1.5) -- (4.2, 2.2) -- (4.5, 2.0) -- (4.8, 2.2) -- (4.5, 1.5); % simple ship shape
  \node[anchor=north, font=\small, neutral70] at (4.5, 1.0) {Target};

  % Ambiguity band (where target could actually be)
  % A vertical band at target range showing elevation uncertainty
  \fill[rose, opacity=0.25] (4.3, 1.5) rectangle (4.7, 3.2);
  \draw[line width=1.5pt, rose, densely dotted] (4.3, 1.5) -- (4.3, 3.2);
  \draw[line width=1.5pt, rose, densely dotted] (4.7, 1.5) -- (4.7, 3.2);

  % Vertical ambiguity annotation
  \draw[line width=1pt, rose, -latex] (5.0, 1.5) -- (5.0, 3.2);
  \node[anchor=west, font=\small, rose, text width=2cm, align=left] at (5.1, 2.35) {Elevation Ambiguity};

  % Range annotation
  \draw[line width=1pt, neutral50, -latex] (0.5, 0.1) -- (4.5, 0.1);
  \node[anchor=north, font=\small, neutral50] at (2.5, -0.2) {Range $r$};

  % ============ RIGHT SIDE: CAMERA VIEW ============

  % Title for right section
  \node[anchor=south, font=\bfseries\Large] at (10.5, 9.5) {Camera Image Projection};

  % Horizon line
  \draw[line width=1.5pt, seablue] (8, 4.0) -- (13, 4.0);
  \node[anchor=north, font=\small, seablue] at (13.2, 4.0) {Horizon};

  % Frustum representation (camera field of view)
  \draw[line width=1.5pt, neutral70] (8, 1.5) -- (8, 6.5);
  \draw[line width=1.5pt, neutral70] (13, 1.5) -- (13, 6.5);
  \draw[line width=1.5pt, neutral70] (8, 1.5) -- (13, 1.5);
  \draw[line width=1.5pt, neutral70] (8, 6.5) -- (13, 6.5);
  \node[anchor=east, font=\tiny, neutral70, text width=1.5cm, align=right] at (7.8, 4.0) {Camera Frustum};

  % Ship silhouette at water level (true location)
  \draw[line width=1.5pt, neutral70, fill=neutral70, opacity=0.6] (10.3, 3.8) -- (10.0, 4.5) -- (10.3, 4.2) -- (10.6, 4.5) -- (10.3, 3.8);

  % Vertical projection line (ambiguity band in camera)
  % This represents the entire vertical range of possible elevations
  \fill[rose, opacity=0.25] (10.0, 3.0) rectangle (10.6, 5.5);
  \draw[line width=2pt, rose] (10.3, 3.0) -- (10.3, 5.5);
  \draw[line width=1.5pt, rose, densely dotted] (10.0, 3.0) -- (10.0, 5.5);
  \draw[line width=1.5pt, rose, densely dotted] (10.6, 3.0) -- (10.6, 5.5);

  % Bounding box uncertainty
  \draw[line width=1.5pt, rose, densely dashed] (10.0, 3.0) rectangle (10.6, 5.5);
  \node[anchor=west, font=\small, rose, text width=2cm, align=left] at (10.7, 4.25) {Detection Uncertainty};

  % "Assumed height" annotation
  \draw[line width=1pt, neutral70, -latex] (9.5, 3.9) -- (10.0, 3.9);
  \node[anchor=east, font=\small, neutral70, text width=2cm, align=right] at (9.4, 3.9) {Assumed sea-surface};

  % Upper uncertainty (could be sky/mast)
  \draw[line width=1pt, rose, -latex] (11.0, 4.5) -- (11.0, 5.5);
  \node[anchor=west, font=\tiny, rose, text width=1.5cm, align=left] at (11.1, 5.0) {Unknown elevation};

  % ============ CONNECTING ARROW ============
  \draw[line width=2pt, neutral50, -latex, bend right=30] (7.5, 4.0) to (8.0, 4.0);
  \node[anchor=north, font=\small, neutral50] at (7.75, 3.7) {Projection};

  % ============ LEGEND / EXPLANATION ============
  \node[anchor=north, font=\bfseries, align=center, text width=12cm] at (6.5, 0.3) {%
    The 2D radar beam (elevation only) creates a vertical ambiguity band at the target's range and azimuth.
    Without elevation information, the projected detection becomes a vertical line in the camera image.%
  };

\end{tikzpicture}
\caption{The three-dimensional ambiguity problem for 2D marine radar. Without elevation measurement, a radar detection at range $r$ and azimuth $\theta$ maps to a vertical line in the camera image rather than a single point. The assumed sea-surface height constrains but does not eliminate vertical uncertainty.}
\label{fig:3d_ambiguity}
\end{figure}
